\section{Introduction}\label{sec:introduction}

For a nonzero unital $C^*$-algebra $A$ and a free ultrafilter $\omega$ on
$\N$, write
\[
 A_\omega=\ell^\infty(\N,A)\big/
 \{(a_n):\lim_{n\to\omega}\|a_n\|=0\},
 \qquad F_\omega(A)=A_\omega\cap A',
\]
where $A$ is identified with its diagonal image in $A_\omega$. A
\emph{character} is a nonzero $*$-homomorphism to $\C$; on a unital algebra
it is automatically unital. We write $\Z$ for the Jiang--Su algebra
introduced by Jiang and Su \cite{JS}. An algebra is \emph{$\Z$-stable}
if it is isomorphic to its minimal tensor product with $\Z$.

Tensorial absorption can be recognized through central copies of the
absorbed algebra. The theory of strongly self-absorbing algebras developed
by Toms and Winter makes this principle precise \cite{TW}. In the case
of $\Z$, Kirchberg and R\o rdam proved that a unital separable algebra
is $\Z$-stable exactly when its norm central-sequence algebra contains,
with the same unit, a separable subhomogeneous algebra without characters
\cite[Theorem~2.5]{KR}. Here subhomogeneity means a common finite bound
on the dimensions of irreducible representations. The absence of
characters on that subalgebra plainly forces the same absence on the
whole central-sequence algebra. The question is whether this elementary
obstruction is already sufficient.

Kirchberg and R\o rdam asked whether the absence of characters on
$F_\omega(A)$ characterizes Jiang--Su absorption among unital separable
$C^*$-algebras \cite[Question~3.1]{KR}. Their work relates this question to
tensor powers of algebras without characters
\cite[Question~3.7 and Proposition~3.9]{KR}. The character question also
appears as Problem~LXXX in the recent collection of
Schafhauser, Tikuisis, and White \cite[Section~22]{STW}. We prove the
proposed characterization.

\begin{theorem}\label{thm:main}
Let $A$ be a nonzero unital separable complex $C^*$-algebra, and let
$\omega$ be any free ultrafilter on $\N$. Then
\[
 F_\omega(A)\text{ has no characters}
 \quad\Longleftrightarrow\quad A\cong A\otmin\Z.
\]
\end{theorem}

The hypothesis uses the norm central-sequence algebra. The theorem applies
to all such $A$, with no nuclearity, simplicity, trace, dimension, or
comparison assumption. Its main ingredient is the following finite
construction, in which separability is unnecessary.

\begin{theorem}\label{thm:finite}
Let $D$ be a nonzero unital complex $C^*$-algebra without characters. There
is an integer $m\ge1$ and a unital $*$-homomorphism
\[
 I(2,3)\longrightarrow D^{\otimes_{\max}m},
\]
where
\[
I(2,3)=\{f\in C([0,1],M_2\otimes M_3):
 f(0)\in M_2\otimes1_3,\ f(1)\in1_2\otimes M_3\}.
\]
\end{theorem}

The exponent depends on $D$. The conclusion is a unital homomorphism,
not a claim that this finite tensor power itself absorbs $\Z$.
Its image is subhomogeneous and characterless, even if the homomorphism
has a kernel. These are precisely the properties needed for the
central-sequence application. They also give the following consequence.

\begin{corollary}[Infinite minimal tensor powers]\label{cor:infinite-tensor}
Let $D$ be a nonzero unital separable complex $C^*$-algebra without
characters, and write
\[
 B=D^{\otimes_{\min}\infty}
 :=\varinjlim_{n\ge1}
 \bigl(D^{\otimes_{\min}n},\,x\longmapsto x\otimes1_D\bigr).
\]
Then $B\cong B\otmin\Z$, and $B$ contains a unital copy of $\Z$.
\end{corollary}

Separability is needed for this application of Dadarlat--Toms, but neither
nuclearity nor simplicity is assumed. The absorption conclusion concerns
the infinite minimal tensor power, not a finite power.

\subsection{Earlier results and the remaining construction}

Dadarlat and Toms connected Jiang--Su absorption with infinite tensor
powers \cite{DT}. Their Theorem~1.1 shows that the infinite minimal
tensor power of a unital separable algebra is $\Z$-stable whenever
it contains a same-unit subhomogeneous characterless subalgebra. The
construction in their Theorem~6.2 gives, for each fixed prime
dimension-drop algebra, a unital map into any separable recursive
subhomogeneous algebra of finite topological dimension whose minimum
matrix size is sufficiently large relative to that dimension.
Our finite theorem begins only with absence of characters on $D$ and
uses maximal tensor products, so it must construct the required
subhomogeneous image without such a structural hypothesis.

A positive element is \emph{full} if the closed two-sided ideal it
generates is the whole ambient algebra. Several regularity criteria
explain why finding the subhomogeneous image is useful.
Robert and Tikuisis showed that two full orthogonal positive elements in
the central-sequence algebra suffice for absorption when the original
separable algebra has finite nuclear dimension \cite[Theorem~1.2]{RT}.
More recently, Perera, Thiel, and Vilalta characterized $\Z$-stability
of separable algebras by purity of the central-sequence algebra
\cite[Theorem~3.6]{PTV}.
Purity imposes comparison and divisibility conditions on its Cuntz
semigroup. The present argument assumes neither finite nuclear dimension
nor these comparison and divisibility properties.

The companion \cite[Theorem~1.4]{Fixed231} proves that a finite maximal tensor power
of every nonzero unital characterless algebra contains two full
orthogonal positive elements, without separability or nuclearity.
Its Gaussian and pure-state-excision proof is different from the
representation-label argument developed below. We also use the
companion's commuting-copy lemma \cite[Lemma~9.3]{Fixed231}, with the
primary reindexing provenance in
\cite[Proposition~2.2 and Corollary~2.3]{KR}.

The bridge from full orthogonal elements to a homomorphism from the matrix
cone $C_0((0,1])\otimes M_2$ that sends $t\mapsto t e_{11}$ to a full
element already appears in the proof of
\cite[Proposition~6.1]{KR}. We give a complete
construction using representation-stable tensor labels. A matrix-cone
map can equivalently be viewed as a completely positive contractive
order-zero map on a matrix algebra, meaning one that preserves
orthogonality of positive elements,
as in Winter and Zacharias \cite[Corollary~4.1]{WZ}; below we construct
the cone homomorphism directly. Fullness of its first matrix corner is
not by itself the dimension-drop criterion. The additional task is to
produce two \emph{exactly equivalent} orthogonal positive contractions,
of the form $v^*v$ and $vv^*$ for a single element $v$, and compare
their unit defect with a positive cutoff of one of them.
This is the input to the stable-rank-free implication of
R\o rdam and Winter \cite[Proposition~5.1, (ii)$\Rightarrow$(iv)]{RW},
whose needed form is proved explicitly in Section~\ref{sec:rw}.

\subsection{Proof overview}

We give the finite construction in two stages. First, a finite-dimensional
space of selfadjoints detects noncommutativity in every representation of
$D$. An incidence argument supplies tensor labels that remain linearly
independent after arbitrary rank-at-least-two maps on the individual legs.
These labels yield two selfadjoints that do not commute in any irreducible
representation of a suitable tensor power. A double commutator then has
full positive and negative parts, and a second use of the labels assembles
finitely many bridges between them into one full square-zero element.

Second, the square-zero element gives a matrix-cone homomorphism with a full
first corner. A
continuous simplex of embeddings of $M_2$ organizes this cone across a
finite tensor power. Two orthogonal equivalent positive elements leave a
defect concentrated near the origin of a cube. Disjoint regions near its
coordinate axes support elements that absorb this defect after a further
finite amplification. The resulting comparison is precisely the one
required by the R\o rdam--Winter dimension-drop criterion. This comparison
is the step that turns full orthogonal elements into a unital $I(2,3)$ map.

Section~\ref{sec:preliminaries} records the comparison conventions and
structural criteria, with the dimension-drop conversion proved in
Section~\ref{sec:rw}. Section~\ref{sec:tensor} constructs this cone homomorphism,
and Section~\ref{sec:dimension-drop} proves the finite comparison.
Section~\ref{sec:application} deduces both theorems and
Corollary~\ref{cor:infinite-tensor}. All tensor exponents and comparison
constants are allowed to depend on the algebra $D$.
